\documentclass[12pt,a4paper]{article}
\usepackage[utf8]{inputenc}
\usepackage[spanish]{babel}
\usepackage{geometry}
\usepackage{booktabs}
\usepackage{tabularx}
\usepackage{float}

\geometry{margin=2.5cm}

\title{Reglas de Negocio Seleccionadas}
\author{}
\date{}

\begin{document}

\maketitle

\section{Reglas de Negocio Seleccionadas}

Las siguientes reglas de negocio son directrices específicas que definen o restringen ciertos aspectos del funcionamiento del sistema, organizadas según las reglas RN07, RN08, RN10 y RN14.

\begin{table}[H]
\centering
\renewcommand{\arraystretch}{1.2}
\begin{tabularx}{\textwidth}{@{} l l X @{}}
\toprule
\textbf{Código} & \textbf{Nombre} & \textbf{Descripción} \\
\midrule

RN07 & Idioma de procesamiento OCR &
El sistema de OCR únicamente procesa tickets en idioma español. Tickets en otros idiomas requieren registro manual. \\

RN08 & Antigüedad recomendada de tickets &
Se recomienda fotografiar tickets al momento de la compra. Tickets con más de 30 días de antigüedad pueden tener menor calidad de escaneo. \\

RN10 & Validación manual obligatoria &
Cuando el OCR presenta discrepancias o baja confianza en la extracción, se solicita validación manual obligatoria del usuario. \\

RN14 & Categorización automática &
Los gastos se clasifican automáticamente según palabras clave y patrones de coincidencia con categorías predefinidas. \\

\bottomrule
\end{tabularx}
\caption{Reglas de Negocio Seleccionadas (RN07, RN08, RN10, RN14)}
\end{table}

\end{document}